\documentclass[10pt,a4paper]{amsart}
\usepackage[margin=19mm]{geometry}
\usepackage{amsmath,amssymb,mathtools}
\usepackage[scr=rsfs,cal=euler]{mathalfa}
\usepackage{xfrac}
\usepackage[unicode,hidelinks,bookmarks=false]{hyperref}
\newcommand{\ie}{i.e.\ }
\newcommand{\cf}{cf.\ }
\newcommand{\resp}{respectively\ }
\newcommand{\Subsection}{\subsection}
\providecommand{\nobreakdash}{\nobreak\hbox{-}}
\setlength{\emergencystretch}{2em}
\multlinegap0pt
\usepackage{bm}
\usepackage{bbm}
\usepackage{dsfont}
\usepackage{xspace}
\usepackage{cancel}
\usepackage{mathtools}
\usepackage[shortlabels]{enumitem}
\usepackage{amsthm}
\makeatletter
\@ifundefined{theorem}{\newtheorem{theorem}{Theorem}}{}
\makeatother
\title[Non-Hausdorff Incidence Completions of Finite Coverings: Mon]{Non-Hausdorff Incidence Completions of Finite Coverings: Monodromy, Transport, and Defect Posets}
\author{Abhiram Sripat}
\date{Source: arXiv:2506.00031v6 (CC BY 4.0)}
\begin{document}
\maketitle

\noindent\emph{Source and licence.} Non-Hausdorff Incidence Completions of Finite Coverings: Monodromy, Transport, and Defect Posets. Version arXiv:2506.00031v6 (CC BY 4.0), distributed under the Creative Commons Attribution 4.0 International licence, as recorded in the arXiv licence metadata for this exact version and shown on the abstract page https://arxiv.org/abs/2506.00031v6. Full rights notice: licenses/arxiv-2506.00031-CC-BY-4.0.txt.\par\smallskip

\noindent\textit{Editorial scope.} Counted unit: the Theorem whose source label is \texttt{thm:incidence-lift} in the article (source0.tex lines 1083--1122, source label \texttt{thm:incidence-lift}), delivered together with the article's own complete original proof of it (source0.tex lines 1124--1318, one \texttt{proof} environment of 195 lines). No mathematics is written, repaired or re-derived: statement and proof are the article's own text line for line. Required context retained uncounted: none. Every definition, display and label of those lines is retained unchanged. Omitted from this package and not redistributed: 1--1082, 1123--1123, 1319--6627. Nothing inside a retained excerpt is omitted or altered: every retained body is delivered exactly as the article's lines are written, so no omission receipt of any kind accompanies this package. Citations: the retained excerpts carry no citation command at all, so no bibliography entry of the article is transcribed and no reference is redistributed. Reference adaptation: none was needed and none was made; every reference inside the delivered unit resolves inside it, so no unresolved reference or citation is produced. Printed slips kept literal: no printed word, symbol, hypothesis or number of the counted statement or of its proof was altered. Layout and class: the article's own class is \texttt{article} with packages including inputenc, fontenc, amsmath, amssymb, amsthm, mathtools, xcolor, hyperref, listings, cleveref, tikz, centernot, amsfonts, mathrsfs; the wrapper is the lane's pinned \texttt{amsart} editorial wrapper at 10pt on A4, which restates the theorem-like environments the retained text needs and adds the article's own macro definitions that the retained text needs (none), with extra wrapper packages (bm, bbm, dsfont, xspace, cancel, mathtools, enumitem). The delivered numbering is the unit's own. Assets: no retained span contains a figure environment, a graphics-inclusion command, a PDF-page inclusion, a table or a TikZ picture; no image file is copied into this package, the package declares no asset dependency and no image file is redistributed with it. Licence: version arXiv:2506.00031v6 of this article is distributed under the Creative Commons Attribution 4.0 International licence, as recorded in the arXiv licence metadata for this exact version and shown on the abstract page https://arxiv.org/abs/2506.00031v6. A full rights notice is delivered at licenses/arxiv-2506.00031-CC-BY-4.0.txt. Characters: every retained excerpt is pure ASCII, so the delivered engine, XeLaTeX with its default Latin Modern text font, reports no missing character.
\begin{theorem}[Incidence lift theorem]
\label{thm:incidence-lift}
Let $\mathcal D$ be an admissible incidence datum and let
\[
\widetilde f^\circ:
Z\setminus D\to Y^\circ
\]
be a fixed regular lift of $f$.

An exceptional colouring
\[
\alpha:D\to\bigsqcup_iF_i
\]
extends $\widetilde f^\circ$ to a continuous lift
\[
\widetilde f:
Z\longrightarrow\overline Y_{\mathcal D}
\]
if and only if the following conditions hold:
\begin{enumerate}[label=(\roman*)]
\item each
\[
\alpha_i:D_i\to F_i
\]
is continuous;
\item for every $z\in D_i$,
\[
\alpha_i(z)
\in
\bigcap_{c\in
\Lambda_{\widetilde f^\circ}(z)}
A_{i,c}.
\]
\end{enumerate}
For the prescribed pair
\[
(\widetilde f^\circ,\alpha)
\]
the extension is unique.
\end{theorem}

\begin{proof}
Suppose first that
\[
\widetilde f:Z\to\overline Y_{\mathcal D}
\]
is a continuous extension.  Its restriction to $D_i$ is continuous
because $F_i$ has its prescribed subspace topology.

Fix
\[
z\in D_i
\]
and let
\[
c\in\Lambda_{\widetilde f^\circ}(z).
\]
Set
\[
a:=\widetilde f(z)\in F_i.
\]
Assume, toward a contradiction, that
\[
a\notin A_{i,c}.
\]
Since $A_{i,c}$ is closed, there exists an open neighbourhood
\[
V\subseteq F_i
\]
of $a$ such that
\[
V\cap A_{i,c}=\varnothing.
\]
Consequently
\[
c\notin C_i(V).
\]

Choose a basic neighbourhood
\[
B_i(V,\mathbf r)
\]
of $a$.  By construction it contains no terminal tail of the component
$c$.

Continuity of $\widetilde f$ at $z$ implies that some neighbourhood
$W$ of $z$ satisfies
\[
\widetilde f(W)
\subseteq
B_i(V,\mathbf r).
\]
But
\[
c\in\Lambda_{\widetilde f^\circ}(z)
\]
means that $W$ contains regular points $w$ with
\[
\widetilde f^\circ(w)\in E_{i,c}.
\]
After shrinking $W$ if necessary, continuity of $f$ forces such points
arbitrarily far down the $c$-tail.  None of these points can belong to
$B_i(V,\mathbf r)$, a contradiction.  Hence
\[
a\in A_{i,c}.
\]
Since this holds for every accumulated end $c$,
\[
a\in
\bigcap_{c\in\Lambda_{\widetilde f^\circ}(z)}
A_{i,c}.
\]

Conversely, assume (i) and (ii).  Define
\[
\widetilde f(z)
=
\begin{cases}
\widetilde f^\circ(z),
&
z\notin D,
\\[1mm]
\alpha_i(z),
&
z\in D_i.
\end{cases}
\]
The map lies over $f$ by construction.  Continuity away from $D$ is
already known.

Fix
\[
z\in D_i
\]
and let
\[
B_i(V,\mathbf r)
\]
be a basic neighbourhood of
\[
\alpha_i(z).
\]

By continuity of $\alpha_i$, there exists a neighbourhood $W_0$ of $z$
in $Z$ such that
\[
\alpha_i(W_0\cap D_i)\subseteq V.
\]
Because $\Sigma$ is finite and $f(z)=z_i$, after shrinking $W_0$ we may also
assume
\[
W_0\cap D_j=\varnothing
\qquad (j\neq i).
\]

Put
\[
C_i^-:=C_i\setminus C_i(V).
\]
If
\[
c\in C_i^-,
\]
then
\[
A_{i,c}\cap V=\varnothing.
\]
Since
\[
\alpha_i(z)\in V,
\]
condition (ii) implies
\[
c\notin\Lambda_{\widetilde f^\circ}(z).
\]
Therefore there exists a neighbourhood $W_c$ of $z$ satisfying
\[
W_c\cap
(\widetilde f^\circ)^{-1}(E_{i,c})
=
\varnothing.
\]
Because $C_i^-$ is finite, the intersection
\[
W_1
:=
W_0\cap
\bigcap_{c\in C_i^-}W_c
\]
is still a neighbourhood of $z$.

For the remaining ends $c\in C_i(V)$, when $C_i(V)\neq\varnothing$ set
\[
r_0:=\min\left\{r_i^*,\min_{c\in C_i(V)}r_c\right\}.
\]
Continuity of $f$ gives a neighbourhood $W_2$ of $z$ such that every
regular point of $W_2$ maps into
\[
U_i(r_0)^\times.
\]

Hence, for
\[
w\in
W_1\cap W_2\setminus D,
\]
the lift $\widetilde f^\circ(w)$ lies in a peripheral component belonging
to $C_i(V)$ and sufficiently far down its tail to lie in
\[
B_i(V,\mathbf r).
\]
Defect points sufficiently near $z$ are mapped by $\alpha_i$ into $V$.
Thus
\[
\widetilde f(W_1\cap W_2)
\subseteq
B_i(V,\mathbf r),
\]
proving continuity at $z$.

If
\[
C_i(V)=\varnothing,
\]
then $C_i^-=C_i$.  Intersect $W_1$ with a neighbourhood $W_2$ on which
\[
f(W_2)\subseteq U_i(r_i^*).
\]
Every regular point of $W_2$ would then lift into one of the representatives
$E_{i,c}$, while $W_1$ excludes all such representatives.  Hence
$W_1\cap W_2$ contains no regular points, and continuity reduces to that of
$\alpha_i$.

The extension is unique once its values on $Z\setminus D$ and on $D$ are
prescribed.
\end{proof}

\end{document}
